% ============================================================================
% HodgeCY v4 revision inserts, keyed to the v3 manuscript.
% Contents:
%   [A] Replacement abstract
%   [B] Patches to §1.5 (contributions) and §1.6 (nonclaims)
%   [C] NEW §4.6: Verified node schemes (Theorem), replacement Table 4
%   [D] Replacement core of §6: critical degree fixed; universal defect
%       theorem; computed defect theorem; postulation-equivalence;
%       strict-refinement corollary; replacement Table 6
%   [E] NEW §6.7: fusion and separation conjectures, sharpened
%   [F] Patches to §7 (implementation), §8.1–8.2, §9 (conclusion)
%   [G] New bibliography entries
%   [H] Figure suggestion (p4-collinearity graphs)
% Conventions follow v3: nodes = ordinary double points; Sigma = node set;
% I_Sigma its homogeneous ideal in Q[x,y,z,t]; h^0(O_{P^3}(8)) = 165.
% ============================================================================

% ----------------------------------------------------------------------------
% [A] REPLACEMENT ABSTRACT
% ----------------------------------------------------------------------------
\begin{abstract}
We develop HodgeCY, a computational framework for studying Hodge atom
profiles of singular and nodal Calabi--Yau threefolds. The first testbed is
the Cynk--Meyer dataset of double octic Calabi--Yau threefolds arising from
eight-plane arrangements, whose control pair $84$/$84^{a}$ has identical
singularity-count vectors and Hodge numbers but different modular forms. We
construct a smoothing bridge from arrangement geometry to finite-node
conifold geometry: for $A = P_1 \cdots P_8$ we study branch octics
$F_\varepsilon = A + \varepsilon Q^2$, and we verify, for the explicit
quartic $Q_0 = x^4 + 2y^4 + 3z^4 + 5t^4 + xyzt$ and $\varepsilon = 1$, that
the resulting branch surfaces for $84$ and $84^{a}$ have exactly $112$
ordinary nodes and no other singularities, so both smoothing bridges are
genuine $112$-node conifold double covers. The paper proves two separations
and one coincidence. First, first-order plane-node incidence is universal
(four copies of the $K_8$ incidence matrix) and separates nothing. Second,
the intersection-lattice concurrency profiles of $84$ and $84^{a}$ are
non-isomorphic, detected by double-line point-type multiplicities and
$p_4$-collinearity graphs. Third, fixing the critical degree $8 = 3k-4$
from classical double-solid theory, we prove that the classical nodal
defects of the two verified $112$-node configurations are \emph{equal},
$\delta(84) = \delta(84^{a}) = 7$, and indeed that the two node schemes have
identical Hilbert functions in every degree. The value $7$ is explained
structurally: we prove a universal lower bound $\delta \geq 7$ for every
$l_3 = 0$ eight-plane smoothing bridge, generated by the seven syzygies
among the plane-product forms $A/P_i$, so the classical defect is
lattice-blind in this family. Consequently the HodgeCY concurrency layer
strictly refines the classical defect on the Cynk--Meyer control pair, and
any separation of the smoothed Hodge atom spectra must occur beyond the
postulation of the node scheme. The remaining operator-side computations are
recorded as explicit validation gates in the public HodgeCY repository.
\end{abstract}

% ----------------------------------------------------------------------------
% [B] PATCHES TO §1.5 AND §1.6
% ----------------------------------------------------------------------------
% §1.5: REPLACE the third, sixth and seventh contribution paragraphs with the
% following (first, second, fourth, fifth stand as in v3, renumbered).

%% (replaces "Third, ...")
Third, we construct and analyze the smoothing bridge
$F_\varepsilon = P_1 \cdots P_8 + \varepsilon Q^2$, and we complete its
node-scheme verification. For the explicit quartic
$Q_0 = x^4 + 2y^4 + 3z^4 + 5t^4 + xyzt$ and $\varepsilon = 1$, the
genericity conditions (G1)--(G2) hold by exact rational computation, and the
singular locus of the branch octic is a reduced zero-dimensional scheme of
length $112$ equal to the union of the $28$ double-line quadruples, with
every singular point an ordinary node (Theorem~\ref{thm:verified-nodes}).
Both arrangements therefore give verified, not merely expected, $112$-node
conifold double covers.

%% (replaces "Sixth, ...")
Sixth, we fix the critical degree of the classical nodal defect for double
octics from the Clemens--Cynk double-solid literature, namely degree
$8 = 3k - 4$ with $2k = 8$, and we compute the defect of both verified
node schemes: $\delta(84) = \delta(84^{a}) = 7$
(Theorem~\ref{thm:defect-computed}). The two node schemes in fact have
identical Hilbert functions in every degree
(Theorem~\ref{thm:postulation}), so no invariant factoring through the
postulation of the node scheme separates the two smoothing bridges.

%% (insert after "Sixth" as a new contribution; renumber "Seventh" of v3 to
%% "Eighth")
Seventh, we explain the common defect value structurally. For every
eight-plane arrangement with $l_3 = 0$ and every smoothing bridge satisfying
(G1)--(G3), the eight degree-seven plane products $A/P_i$ vanish on all
$28$ double lines, their linear multiples span a $25$-dimensional space cut
by exactly seven syzygies, and this space meets $Q \cdot H^0(\mathcal{O}(4))$
trivially. This yields the universal lower bound $\delta \geq 7$
(Theorem~\ref{thm:universal-defect}), realized with equality by $84$ and
$84^{a}$. The universal part of the defect is therefore lattice-blind, in
precise analogy with the universality of first-order plane-node incidence
(Proposition~5.1). It follows that the concurrency layer strictly refines
the classical defect on the control pair
(Corollary~\ref{cor:strict-refinement}): this resolves the three-outcome
fork of Section~6.3 in its second, strongest branch.

% §1.6 (Scope and nonclaims): REPLACE the second paragraph ("We also do not
% claim that the classical defects ... are equal or different.") with:

The classical defects of the two verified $112$-node configurations are now
computed and equal, $\delta = 7$ for both, and the full Hilbert functions of
the node schemes agree. What we still do not claim is that the smoothed
$84$ and $84^{a}$ models have different Hodge atom spectra. The completed
computations sharpen rather than settle that question: they prove that any
such separation, if it exists, is invisible to node count, first-order
incidence, classical defect, and the entire postulation of the node scheme,
and must therefore be carried by finer relation-sensitive or operator-side
structure. This is formalized in the conjectures of Section~6.7.

% Also in §1.6, DELETE the sentence in the final paragraph reading "For
% explicit examples, one must still verify the singular locus, node count,
% and local quadratic rank by computer algebra." and REPLACE with:
For the pair $84$/$84^{a}$ this verification is carried out in
Section~4.6; for other arrangements it remains a per-example computational
gate.

% ----------------------------------------------------------------------------
% [C] NEW §4.6 — insert at the end of Section 4
% ----------------------------------------------------------------------------
\subsection{Verified node schemes for arrangements 84 and $84^{a}$}
\label{sec:verified-nodes}

The genericity analysis of Sections~4.1--4.4 gives the local mechanism; this
subsection removes the word ``expected'' for the pair $84$/$84^{a}$ by
verifying an explicit smoothing datum. Fix
\[
  Q_0 = x^4 + 2y^4 + 3z^4 + 5t^4 + xyzt,
  \qquad \varepsilon = 1,
  \qquad F = P_1 \cdots P_8 + Q_0^2 ,
\]
and write $D = \{F = 0\}$, $\Sigma = \operatorname{Sing}(D)$, and
$J = (\partial_x F, \partial_y F, \partial_z F, \partial_t F)$ for the
Jacobian ideal. Let
\[
  I_\Sigma \;=\; \bigcap_{1 \le i < j \le 8}
     \bigl( P_i,\, P_j,\, Q_0 \bigr)
\]
be the ideal of the $28$ double-line quadruples
$\bigcup_{i<j} \bigl( L_{ij} \cap \{Q_0 = 0\} \bigr)$.

We first record two elementary lemmas; the first replaces a Gr\"obner
containment test by a one-line argument, and the second reduces nodality to
reducedness.

\begin{lemma}[Structural containment]\label{lem:containment}
For any eight-plane arrangement and any $(Q, \varepsilon)$, the Jacobian
ideal of $F_\varepsilon = A + \varepsilon Q^2$ satisfies
$J \subseteq I_\Sigma$, where $I_\Sigma$ is the ideal of
$\bigcup_{i<j} (L_{ij} \cap \{Q = 0\})$.
\end{lemma}

\begin{proof}
Let $p \in L_{ij}$ with $Q(p) = 0$. Each partial derivative of $A$ is a sum
of products of seven of the eight plane forms; every summand omits at most
one of $P_i, P_j$ and hence contains the other, so all partials of $A$
vanish at $p$. The partials of $\varepsilon Q^2$ equal
$2\varepsilon Q \,\partial Q$ and vanish at $p$ because $Q(p) = 0$.
\end{proof}

\begin{lemma}[Reducedness certifies nodality]\label{lem:reduced-nodal}
Let $p$ be an isolated singular point of a surface $D \subset \mathbb{P}^3$.
The Jacobian scheme of $D$ is reduced at $p$ if and only if $p$ is an
ordinary node.
\end{lemma}

\begin{proof}
In an affine chart with local equation $f$, the point is an $A_1$ if and
only if the Hessian of $f$ at $p$ has rank $3$, which holds if and only if
the linear parts of $\partial_1 f, \partial_2 f, \partial_3 f$ span the
cotangent space, which holds if and only if the local Jacobian ideal is the
maximal ideal, i.e.\ the Jacobian scheme is reduced at $p$.
\end{proof}

\begin{theorem}[Verified $112$-node schemes]\label{thm:verified-nodes}
Let $A$ be the defining octic of arrangement $84$ or $84^{a}$, and let
$F = A + Q_0^2$ as above. Then:
\begin{enumerate}
\item[(i)] $(Q_0, \varepsilon = 1)$ satisfies the genericity conditions
  (G1)--(G3);
\item[(ii)] the singular scheme of $D = \{F = 0\}$ is reduced,
  zero-dimensional of length $112$, and equals
  $\bigcup_{i<j} \bigl( L_{ij} \cap \{Q_0 = 0\} \bigr)$, four points on
  each of the $28$ double lines;
\item[(iii)] every singular point of $D$ is an ordinary node, and the
  double cover $w^2 = F$ is a $112$-node conifold double cover.
\end{enumerate}
\end{theorem}

\begin{proof}
(i) Conditions (G1) and (G2) are finite statements verified in exact
rational arithmetic: $Q_0$ is nonzero at each of the $26$ multiple points
($16$ triple, $10$ quadruple, matching the Cynk--Meyer counts), and the
restriction of $Q_0$ to each of the $28$ double lines is a squarefree
binary quartic. (G3) is the content of (ii).

(ii) By (G2) each quadruple $L_{ij} \cap \{Q_0 = 0\}$ consists of four
distinct reduced points, and by (G1) the quadruples on distinct lines are
pairwise disjoint, since two double lines meet only at multiple points of
the arrangement, which $Q_0$ avoids. Hence $I_\Sigma$ is a radical ideal
of a reduced scheme of length exactly $112$; this requires no Gr\"obner
computation. By Lemma~\ref{lem:containment}, $J \subseteq I_\Sigma$, so the
projective scheme $V(I_\Sigma)$ is a closed subscheme of the projective
scheme defined by the saturation $J^{\mathrm{sat}}$. A standard-basis
computation over $\mathbb{Q}$ (by modular standard bases, cross-checked at
the primes $32003$ and $1000003$) shows that $V(J)$ is zero-dimensional of
degree $112$. A closed embedding of zero-dimensional projective schemes of
the same finite length is an equality, so
$J^{\mathrm{sat}} = I_\Sigma$: the singular scheme is reduced of length
$112$ and is exactly the union of quadruples. In particular no singular
point exists outside the double lines, which is (G3).

(iii) The singular scheme is reduced at every point by (ii), so every
singularity of $D$ is an ordinary node by Lemma~\ref{lem:reduced-nodal};
the local statement for the double cover is Proposition~4.1.
\end{proof}

\begin{remark}
The proof structure is a certificate: one symbolic containment
(Lemma~\ref{lem:containment}), one exact radical identification from
(G1)--(G2), one degree computation, and one local criterion
(Lemma~\ref{lem:reduced-nodal}). The only heavy computation is the degree
of the Jacobian scheme, and its output is a single integer, making the
verification straightforwardly reproducible in Macaulay2 or Singular. The
scripts are included in the repository
(\texttt{scripts/verify\_node\_scheme}).
\end{remark}

% REPLACEMENT Table 4 (statuses updated):
\begin{table}[htbp]
\caption{Smoothing-bridge profiles for arrangements $84$ and $84^{a}$ with
the verified datum $(Q_0, \varepsilon) = (x^4+2y^4+3z^4+5t^4+xyzt,\ 1)$.}
\begin{tabular}{ccccl}
\hline
Arrangement & Triple lines $l_3$ & Double lines & Nodes per line & Status \\
\hline
$84$   & $0$ & $28$ & $4$ & Verified $112$-node conifold
                             (Theorem~\ref{thm:verified-nodes}). \\
$84^a$ & $0$ & $28$ & $4$ & Verified $112$-node conifold
                             (Theorem~\ref{thm:verified-nodes}). \\
\hline
\end{tabular}
\end{table}

% ----------------------------------------------------------------------------
% [D] REPLACEMENT CORE OF SECTION 6
% ----------------------------------------------------------------------------
% §6.1: APPEND the following paragraphs after the existing definition of the
% defect layer (and delete the sentence "For this reason, we do not insert a
% guessed critical degree in the present paper." together with the sentence
% following it).

The critical degree for double octics is classical. For a double solid
branched along a nodal surface of degree $2k$ in $\mathbb{P}^3$, the defect
is measured by the failure of the nodes to impose independent conditions on
forms of degree $3k - 4$; for sextic double solids this is the familiar
quintic condition, and for octic double solids ($k = 4$) the critical
degree is $8$ \cite{Clemens83, CynkNodalOctics, Cheltsov, KloostermanDefect}.
Concretely, for a node set $\Sigma$ with $|\Sigma| = m$ we take
\[
  \delta(\Sigma)
  \;=\; m \;-\; \operatorname{HF}_{R/I_\Sigma}(8)
  \;=\; h^0\!\bigl(\mathbb{P}^3, \mathcal{I}_\Sigma(8)\bigr)
        - \bigl( h^0(\mathcal{O}_{\mathbb{P}^3}(8)) - m \bigr),
\]
with $h^0(\mathcal{O}_{\mathbb{P}^3}(8)) = 165$, so that $\delta = 0$ exactly
when the nodes impose independent conditions in degree $8$. Two consistency
checks from the literature: factoriality of a nodal octic double solid is
equivalent to independence of the node conditions in degree $8$
\cite{Cheltsov}, and a nodal double cover branched in degree $2d$ with
positive defect has at least $d(2d-1)$ nodes \cite{KloostermanDefect} ---
$28$ for $d = 4$, so the $112$-node bridges lie well inside the range where
defect phenomena are possible.

% §6.4 of v3 ("Current defect-computation queue"): REPLACE the entire
% subsection with the following two subsections. (The v3 verification plan
% §6.5 should be retitled "Independent confirmation plan" and its seven steps
% kept, with a sentence noting that steps 1–5 and 7 are carried out in
% Sections 4.6 and 6.4–6.5 and that the Macaulay2 templates remain in the
% repository for independent confirmation.)

\subsection{The universal defect of $l_3 = 0$ smoothing bridges}
\label{sec:universal-defect}

The defect computation for the verified bridges turns out to have a
structural core that is independent of the arrangement lattice. Throughout
this subsection $A = P_1 \cdots P_8$ is an arrangement of eight distinct
planes with $l_3 = 0$, and $F_\varepsilon = A + \varepsilon Q^2$ is a
smoothing bridge satisfying (G1)--(G3), with node set $\Sigma$ of length
$112$.

\begin{lemma}\label{lem:plane-products}
For each $i$, the degree-seven form $A/P_i$ vanishes on every double line
$L_{jk}$, and hence lies in $I_\Sigma(7)$.
\end{lemma}

\begin{proof}
If $i \notin \{j,k\}$ then $P_j$ divides $A/P_i$; if $i = j$ then $P_k$
divides $A/P_i$. Either way a factor vanishes on $L_{jk}$. Since
$\Sigma \subset \bigcup L_{jk}$, the form lies in $I_\Sigma$.
\end{proof}

\begin{lemma}\label{lem:syzygies}
Let $W \subseteq H^0(\mathcal{O}_{\mathbb{P}^3}(8))$ be the span of the
products $(A/P_i)\,\ell$ with $\ell$ linear. Then $\dim W = 25$: the kernel
of the assembly map
$\bigoplus_{i=1}^{8} H^0(\mathcal{O}(1)) \to H^0(\mathcal{O}(8))$,
$(\ell_i) \mapsto \sum_i (A/P_i)\ell_i$, is exactly the seven-dimensional
space $\{(c_i P_i) : \sum_i c_i = 0\}$.
\end{lemma}

\begin{proof}
Suppose $\sum_i (A/P_i)\ell_i = 0$. Reducing modulo $P_i$ kills every
summand except the $i$-th, so $P_i$ divides $(A/P_i)\ell_i$; since $P_i$
does not divide $A/P_i$, it divides $\ell_i$, so $\ell_i = c_i P_i$. Then
$\sum_i c_i A = 0$ forces $\sum_i c_i = 0$. Conversely every such tuple is
a syzygy, since $(A/P_i)P_i = A$ for all $i$. Hence
$\dim W = 8 \cdot 4 - 7 = 25$.
\end{proof}

\begin{lemma}\label{lem:transversal-independence}
Under (G2), $\;Q \cdot H^0(\mathcal{O}(4)) \,\cap\, W = 0$.
\end{lemma}

\begin{proof}
Suppose $Qg \in W$ with $g$ a quartic. Every element of $W$ vanishes on
every double line by Lemma~\ref{lem:plane-products}, so $Qg$ vanishes
identically on each $L_{jk}$. By (G2) the restriction $Q|_{L_{jk}}$ is not
identically zero, and $L_{jk}$ is irreducible, so $g|_{L_{jk}} \equiv 0$
for all $28$ lines. Fix $i$: the plane $P_i$ contains the seven distinct
lines $L_{ij}$, $j \neq i$ (distinct because $l_3 = 0$), and the plane
quartic $g|_{P_i}$ vanishes on all seven. A nonzero quartic curve in
$\mathbb{P}^2$ contains at most four lines, so $g|_{P_i} \equiv 0$, i.e.\
$P_i \mid g$, for every $i$. Then $A \mid g$, impossible for a nonzero
quartic; so $g = 0$.
\end{proof}

\begin{theorem}[Universal defect bound]\label{thm:universal-defect}
For every eight-plane arrangement with $l_3 = 0$ and every smoothing bridge
$F_\varepsilon = A + \varepsilon Q^2$ satisfying (G1)--(G3),
\[
  h^0\!\bigl(\mathcal{I}_\Sigma(8)\bigr) \;\geq\;
  \dim \bigl( Q \cdot H^0(\mathcal{O}(4)) \oplus W \bigr)
  \;=\; 35 + 25 \;=\; 60,
  \qquad\text{hence}\qquad
  \delta(\Sigma) \;\geq\; 60 - 53 \;=\; 7 .
\]
The bound is independent of the projective position of the arrangement: the
sixty-dimensional subspace is generated by the transversal quartic and by
the plane products, and the excess $7 = 8 - 1$ counts the syzygies
$(A/P_i)P_i = A$ among the plane products.
\end{theorem}

\begin{proof}
$Q$ vanishes on $\Sigma$, so $Q \cdot H^0(\mathcal{O}(4)) \subseteq
I_\Sigma(8)$, of dimension $35$ since $Q \neq 0$. $W \subseteq I_\Sigma(8)$
by Lemma~\ref{lem:plane-products}, of dimension $25$ by
Lemma~\ref{lem:syzygies}, and the sum is direct by
Lemma~\ref{lem:transversal-independence}. The expected value of
$h^0(\mathcal{I}_\Sigma(8))$ for $112$ independent points is
$165 - 112 = 53$.
\end{proof}

\begin{remark}\label{rem:defect-lattice-blind}
Theorem~\ref{thm:universal-defect} is the defect-layer analogue of
Proposition~5.1. First-order plane-node incidence is a universal object,
and now the guaranteed part of the classical defect is a universal number:
both are functions of ``eight planes pairwise in general position'' and see
nothing of the intersection lattice. Whether the defect \emph{exceeds} the
universal floor is, a priori, where lattice information could enter; the
next subsection shows that for $84$ and $84^{a}$ it does not.
\end{remark}

\subsection{The computed defect of the verified bridges}
\label{sec:computed-defect}

\begin{theorem}[Defect computation]\label{thm:defect-computed}
For the verified node schemes of Theorem~\ref{thm:verified-nodes},
\[
  \delta(84) \;=\; \delta(84^{a}) \;=\; 7 ,
  \qquad
  h^0\!\bigl(\mathcal{I}_\Sigma(8)\bigr) = 60
  \quad\text{for both arrangements},
\]
and the sixty-dimensional space of degree-eight forms through the nodes is
exactly $Q_0 \cdot H^0(\mathcal{O}(4)) \oplus W$.
\end{theorem}

\begin{proof}
The lower bound $h^0(\mathcal{I}_\Sigma(8)) \geq 60$ is
Theorem~\ref{thm:universal-defect}; for the specific datum it is also
confirmed by an exact rank computation over $\mathbb{Q}$: the $67$ spanning
forms ($35$ products $Q_0 m$ with $m$ a degree-four monomial and $32$
products $(A/P_i)\ell$) have coefficient rank exactly $60$ in the
$165$-dimensional space $H^0(\mathcal{O}(8))$.

For the upper bound, let $p \in \{32003,\, 1000003\}$ and let
$I_{\Sigma,p}$ denote the ideal constructed over $\mathbb{F}_p$ by the same
intersection recipe. In each degree $d$, the graded piece of $I_\Sigma$
over $\mathbb{Q}$ is the kernel of an integer coefficient matrix whose rank
can only drop modulo $p$, and reduction modulo $p$ of an intersection of
ideals is contained in the intersection of the reductions; hence
$\dim_{\mathbb{F}_p} (I_{\Sigma,p})_d \geq \dim_{\mathbb{Q}}
(I_\Sigma)_d$ for every $d$. A standard-basis computation over
$\mathbb{F}_p$ gives $\operatorname{HF}_{R/I_{\Sigma,p}}(8) = 105$, i.e.\
$\dim (I_{\Sigma,p})_8 = 60$, at both primes and for both arrangements.
Therefore $\dim_{\mathbb{Q}} (I_\Sigma)_8 \leq 60$, and equality holds. The
defect is $\delta = 112 - 105 = 7$, and the degree-eight graded piece
coincides with the structural space by dimension count.
\end{proof}

\begin{theorem}[Postulation equivalence]\label{thm:postulation}
The verified node schemes of $84$ and $84^{a}$ have identical Hilbert
functions: for both arrangements,
\[
  \operatorname{HF}_{R/I_\Sigma}(d)
  \;=\; 4,\ 10,\ 20,\ 34,\ 52,\ 74,\ 92,\ 105,\ 112
  \qquad (d = 1, \dots, 9),
\]
with $\operatorname{HF}(d) = 112$ for all $d \geq 9$. In particular the
defect in every degree, not only the critical degree, agrees between the
two configurations.
\end{theorem}

\begin{proof}
The same sandwich as in Theorem~\ref{thm:defect-computed} applies
degreewise. In degrees $d \leq 3$ the finite-field computation gives
$\dim (I_{\Sigma,p})_d = 0$, so $(I_\Sigma)_d = 0$ over $\mathbb{Q}$. In
degrees $4 \leq d \leq 7$ the structural subspaces
$Q_0 \cdot H^0(\mathcal{O}(d-4))$ for $d \leq 6$, and
$Q_0 \cdot H^0(\mathcal{O}(3)) \oplus \operatorname{span}\{A/P_i\}$ of
dimension $20 + 8 = 28$ for $d = 7$ (directness and the eight-dimensionality
of the plane-product span follow as in
Lemmas~\ref{lem:syzygies}--\ref{lem:transversal-independence}), give lower
bounds $1, 4, 10, 28$ that meet the finite-field upper bounds
$35 - 34,\ 56 - 52,\ 84 - 74,\ 120 - 92$. Degree $8$ is
Theorem~\ref{thm:defect-computed}. For $d \geq 9$ the finite-field value
$\operatorname{HF}_p(d) = 112$ forces
$\operatorname{HF}_{\mathbb{Q}}(d) = 112$, since
$\operatorname{HF}_{\mathbb{Q}}(d) \geq \operatorname{HF}_p(d)$ by the
semicontinuity inequality and $\operatorname{HF}_{\mathbb{Q}}(d) \leq 112$
for a length-$112$ scheme.
\end{proof}

\begin{corollary}[Strict refinement of the classical defect]
\label{cor:strict-refinement}
The smoothing bridges of arrangements $84$ and $84^{a}$ have the same node
count, the same universal first-order plane-node incidence, the same
classical defect $\delta = 7$, and the same Hilbert function in every
degree, while their source arrangements have non-isomorphic concurrency
profiles (Theorem~5.4). Hence, on the Cynk--Meyer control pair, the
HodgeCY concurrency layer strictly refines every invariant of the node
scheme that factors through its postulation --- in particular the classical
nodal defect. This realizes the second branch of the three-outcome fork of
Section~6.3.
\end{corollary}

% REPLACEMENT Table 6 (queue -> results):
\begin{table}[htbp]
\caption{Defect computation for the verified smoothing bridges. Critical
degree $8 = 3k-4$ fixed from double-solid theory
\cite{Clemens83, CynkNodalOctics, Cheltsov, KloostermanDefect}; values
proved over $\mathbb{Q}$ by the structural lower bound of
Theorem~\ref{thm:universal-defect} together with finite-field upper bounds
at $p = 32003$ and $p = 1000003$.}
\begin{tabular}{cccccl}
\hline
Arrangement & Nodes & Critical degree & $\operatorname{HF}_{R/I}(8)$ &
Defect $\delta$ & Status \\
\hline
$84$   & $112$ & $8$ & $105$ & $7$ & Computed
  (Theorem~\ref{thm:defect-computed}). \\
$84^a$ & $112$ & $8$ & $105$ & $7$ & Computed
  (Theorem~\ref{thm:defect-computed}). \\
\hline
\end{tabular}
\end{table}

% ----------------------------------------------------------------------------
% [E] NEW §6.7 — insert after the operator-route subsection §6.6
% ----------------------------------------------------------------------------
\subsection{The fusion and separation conjectures}
\label{sec:fusion-separation}

The completed computations organize the remaining atom-level questions into
two precise statements.

\begin{conjecture}[Fusion]\label{conj:fusion}
For an eight-plane arrangement $A$ with $l_3 = 0$, the Hodge atom spectrum
of the smoothed conifold double cover $X_\varepsilon$ is independent of the
generic choice of $(Q, \varepsilon)$ satisfying (G1)--(G3), and is
determined by the intersection-lattice concurrency profile of $A$.
\end{conjecture}

\begin{conjecture}[Separation]\label{conj:separation}
The verified $112$-node conifold double covers $X_\varepsilon(84)$ and
$X_\varepsilon(84^{a})$ have different Hodge atom spectra.
\end{conjecture}

The fusion statement is the well-definedness half: it asserts that the
smoothing bridge is a genuine transfer map from arrangement lattices to
atom spectra. The separation statement asserts that this transfer is not
blind to the lattice on the control pair. The results of this section
sharpen the separation conjecture considerably. By
Theorem~\ref{thm:postulation}, the two node schemes are
postulation-equivalent, so a separation cannot be detected by node count,
incidence, classical defect, or any Hilbert-function invariant of
$\Sigma$; and by Theorem~\ref{thm:universal-defect} the guaranteed part of
the defect is lattice-blind by structural reasons, not by accident.
Whatever distinguishes the two atom spectra, if
Conjecture~\ref{conj:separation} holds, must be carried by the finer
relation geometry of the node configuration --- the block structure of the
$28$ quadruples, the concurrency environment of Theorem~5.4 --- or by
operator-side monodromy along the lines of Section~6.6. Conversely, a
disproof of Conjecture~\ref{conj:separation} would be a rigidity theorem:
the concurrency distinction of Theorem~5.4 would be certified to die in the
smoothing, and the separation hunt would move to the $(3,51)$ cluster of
Section~8.3.


% ----------------------------------------------------------------------------
% [F] PATCHES TO §7, §8, §9
% ----------------------------------------------------------------------------
% §7.2: APPEND after the paragraph on the concurrency separation:

The node-scheme and defect computations of Sections~4.6 and~6.4--6.5 are
implemented in \texttt{scripts/verify\_node\_scheme} and
\texttt{scripts/compute\_defect}. The genericity conditions (G1)--(G2) and
the rank-$60$ identification of the degree-eight graded piece are verified
in exact rational arithmetic (Python/SymPy); the degree of the Jacobian
scheme over $\mathbb{Q}$ is computed by modular standard bases in Singular
and cross-checked at two large primes; the Hilbert functions are computed
in Singular over $\mathbb{F}_{32003}$ and $\mathbb{F}_{1000003}$ and
promoted to $\mathbb{Q}$ by the structural sandwich argument of
Theorem~\ref{thm:defect-computed}. Macaulay2 templates for an independent
confirmation of the same quantities are retained in the repository.

% §7.3, first principle: REPLACE the example sentences ("For example, the
% 112-node count ... have been supplied.") with:

For example, the $112$-node count for arrangements $84$ and $84^{a}$ was
recorded as expected until the explicit verification of
Theorem~\ref{thm:verified-nodes} was completed; it is now recorded as
verified, with the certificate (choice of $Q_0$ and $\varepsilon$, exact
genericity checks, Jacobian-scheme degree, and reducedness) archived
alongside the value. The same convention now applies to the defect layer,
where the recorded value $\delta = 7$ carries its proof route: structural
lower bound over $\mathbb{Q}$, finite-field upper bounds, and the
semicontinuity inequality that closes the sandwich.

% §8.1: APPEND as a new penultimate paragraph:

The fifth conclusion is new in this version: the defect gate is closed, and
it closes in the strongest of the three branches anticipated in
Section~6.3. The classical defects of the two verified $112$-node bridges
are equal, $\delta = 7$, the common value is forced by an
arrangement-universal mechanism (Theorem~\ref{thm:universal-defect}), and
the full Hilbert functions of the two node schemes coincide
(Theorem~\ref{thm:postulation}). The classical numerical theory of nodal
double octics is therefore provably too coarse for the $84$/$84^{a}$
comparison, while the concurrency layer separates the pair. This is the
precise sense in which the HodgeCY profile strictly refines classical
defect theory on canonical examples.

% §8.2: REPLACE items "Fourth" through "Sixth" of the gate list with:

Fourth and fifth, the critical degree and the Hilbert-function defect
computations are complete (Sections~6.4--6.5). Sixth, the comparison is
done: the defects agree, so the decisive open question moves past the
classical layer. The remaining gates are the operator-side import of
step seven and the atom-level formalization of
Conjectures~\ref{conj:fusion} and~\ref{conj:separation}, beginning with a
proof of the fusion statement, for which the natural strategy is a
monodromy argument over the connected parameter space of admissible
$(Q, \varepsilon)$.

% §9: REPLACE the paragraph beginning "The present paper does not claim
% that the smoothed 84 and 84a models have different Hodge atom spectra. It
% also does not claim that their classical defects are equal or different."
% with:

The present paper does not claim that the smoothed $84$ and $84^{a}$ models
have different Hodge atom spectra; that is Conjecture~\ref{conj:separation}.
What it now proves is that the question cannot be settled below the atom
level: the two verified $112$-node conifold double covers have equal
classical defects, $\delta = 7$, with the common value explained by a
universal syzygy mechanism, and their node schemes are
postulation-equivalent in every degree. The separation question is thereby
pushed, with proof, past every classical numerical invariant of the node
scheme, onto exactly the relation-sensitive and operator-side structure
that the HodgeCY profile is designed to record.

% ----------------------------------------------------------------------------
% [G] NEW BIBLIOGRAPHY ENTRIES
% ----------------------------------------------------------------------------
% Add (biblatex/bibtex as appropriate to the manuscript):
%
% @article{CynkNodalOctics,
%   author  = {S{\l}awomir Cynk},
%   title   = {Hodge Numbers of Nodal Double Octics},
%   journal = {Communications in Algebra},
%   volume  = {27}, number = {8}, pages = {4097--4102}, year = {1999}}
%
% @misc{Cheltsov,
%   author = {Ivan Cheltsov},
%   title  = {On Factoriality of Nodal Threefolds},
%   note   = {arXiv:math/0405221; published version: J. Differential Geom.
%             --- verify volume/pages at submission}}
%
% @misc{KloostermanDefect,
%   author = {Remke Kloosterman},
%   title  = {Maximal Families of Nodal Varieties with Defect},
%   note   = {arXiv:1310.0227; published version: Math. Z. --- verify
%             volume/pages at submission}}
%
% (Clemens83 = existing reference [19], "Double Solids".)

% ----------------------------------------------------------------------------
% [H] FIGURE SUGGESTION
% ----------------------------------------------------------------------------
% Replace or supplement Figures 4–5 with the generated two-panel figure
% p4_collinearity_graphs.pdf: the p4-collinearity graphs of 84 and 84a with
% vertices labeled by degree, the unique degree-6 vertex of 84 and the four
% degree-9 vertices of 84a highlighted. Suggested caption:
%
% \caption{The $p_4$-collinearity graphs of arrangements $84$ (left, $39$
% edges, degree sequence $[6, 8^9]$) and $84^{a}$ (right, $42$ edges, degree
% sequence $[8^6, 9^4]$). Vertices are the ten $p_4$ points, labeled by
% degree; two vertices are joined when the corresponding points lie on a
% common double line. Highlighted vertices witness the non-isomorphism of
% Theorem~5.4: the unique degree-$6$ vertex of $84$ and the four degree-$9$
% vertices of $84^{a}$.}
